\documentclass[11pt,a4paper]{article}
\usepackage[utf8]{inputenc}\usepackage[T1]{fontenc}\usepackage[italian,english]{babel}
\usepackage{amsmath,amssymb,amsthm}\usepackage[margin=2.2cm]{geometry}\usepackage{longtable,booktabs,array,calc}
\usepackage{listings,xcolor}\usepackage{hyperref}\usepackage{url}
\providecommand{\tightlist}{\setlength{\itemsep}{0pt}\setlength{\parskip}{0pt}}
\lstset{basicstyle=\ttfamily\scriptsize,breaklines=true,frame=single,captionpos=b,columns=fullflexible,keepspaces=true,inputencoding=utf8,extendedchars=true}
\title{Proof Pursuit --- Problem 4, part 2\\ \large prove, decisioni degli agenti, codice verificato, letteratura}
\author{Team Proof Pursuit (Researcher: T. Tumini; Referee: gabundos; Creative: M. Renzi; gio3r3gio)}
\date{\today}
\begin{document}\maketitle\tableofcontents
\section*{Come leggere questo rapporto}
Per ogni cella: la richiesta ufficiale, lo stato dichiarato dal team (mai ``risolto'' senza revisione umana), la consegna
con la prova, poi la traccia delle decisioni degli agenti (Researcher: famiglia e motivo; Referee: verdetto ed errore
fatale; Creative: quando e perché è intervenuto) e il codice su cui il tentativo poggia con l'esito della riesecuzione
fidata. DIMOSTRATO, CITATO, VERIFICATO SU INTERVALLO FINITO e CONGETTURA restano distinti. L'architettura del sistema
è descritta nel README del repository.
\section{Problem 4 — Disjoint congruence classes and large gcd}
\subsection{Cella 2 (2 punti) — stato: revisionato}
\subparagraph{Parte 2 (C2) --- Four
classes}\label{parte-2-c2-four-classes}

\textbf{Punteggio:} 2 points · \textbf{Valutazione:} Judged

The same question for four classes. Prove the statement for \(k = 4\).

\textbf{Enunciato protetto dato agli agenti (state.json).} \texttt{Prove the statement for k = 4: any four pairwise disjoint classes have a pair with gcd(m\_i,m\_j) \textgreater{}= 4.}

\textbf{Evidenza registrata in STATUS.md.} prova completa (clique per primo, casi su \$

\subsubsection{Consegna (prova)}
\paragraph{Problema 4 --- Parte 2 --- bozza di
consegna}\label{problema-4-parte-2-bozza-di-consegna}

\textbf{Stato dichiarato: RISOLTA (approvata).} Prova completa scritta
dal team; Referee automatico READY\_FOR\_HUMAN (giudice A matematica
PASS, giudice B evidenze PASS); approvazione umana registrata in
\texttt{runs/p4\_c2/approval.json}.

\subparagraph{1. Risultato e ambito}\label{risultato-e-ambito}

Quattro classi disgiunte: l'enunciato della parte è dimostrato
integralmente, nella generalità richiesta (vedi la prova per le ipotesi
esatte). Nessun calcolo al computer è necessario.

\subparagraph{2. Dimostrazione}\label{dimostrazione}

\paragraph{\texorpdfstring{Cell 2: four pairwise disjoint classes have a
pair of moduli with
\(\gcd\ge4\)}{Cell 2: four pairwise disjoint classes have a pair of moduli with \textbackslash gcd\textbackslash ge4}}\label{cell-2-four-pairwise-disjoint-classes-have-a-pair-of-moduli-with-gcdge4}

\textbf{Target.} If \(a_i\ (\mathrm{mod}\ m_i)\), \(i=1,2,3,4\), are
pairwise disjoint congruence classes (\(m_i\ge1\)), then
\(\gcd(m_i,m_j)\ge4\) for some \(i<j\).

\subparagraph{Lemma 0}\label{lemma-0}

Let \(g=\gcd(m,m')\). If \(g\mid a-a'\) then the classes
\(a\ (\mathrm{mod}\ m)\) and \(a'\ (\mathrm{mod}\ m')\) meet;
equivalently, disjoint classes satisfy \(g\nmid a-a'\). \emph{Proof.}
Write \(m=g\mu\), \(m'=g\mu'\) with \(\gcd(\mu,\mu')=1\) and
\(a-a'=g\delta\). The element \(a+g\mu t\) of the first class lies in
the second iff \(g\mu'\mid g(\delta+\mu t)\) iff
\(\mu'\mid\delta+\mu t\), which has the solution
\(t\equiv-\delta\mu^{-1}\pmod{\mu'}\) since \(\mu\) is invertible modulo
\(\mu'\). \(\square\)

\subparagraph{Proof of the cell}\label{proof-of-the-cell}

Write \(g_{ij}=\gcd(m_i,m_j)\) for \(i\ne j\). Suppose, for a
contradiction, that \(g_{ij}\le3\) for all pairs. By Lemma 0,
\(g_{ij}=1\) is impossible for disjoint classes (\(1\mid a_i-a_j\)),
hence \[g_{ij}\in\{2,3\}\quad\text{for all } i\ne j.\tag{1}\] Two
consequences of (1) and Lemma 0, for \(i\ne j\): - (E) if \(2\mid m_i\)
and \(2\mid m_j\), then \(2\mid g_{ij}\), so by (1) \(g_{ij}=2\), and
disjointness gives \(2\nmid a_i-a_j\), i.e.~\(a_i\not\equiv a_j\pmod2\);
- (T) if \(3\mid m_i\) and \(3\mid m_j\), then \(3\mid g_{ij}\), so
\(g_{ij}=3\), and disjointness gives \(a_i\not\equiv a_j\pmod3\).

Let \(E=\{i: 2\mid m_i\}\) and \(T=\{i:3\mid m_i\}\). \emph{Claim 1:
\(|E|\le2\) and \(|T|\le3\).} By (E) the residues \(a_i\bmod2\),
\(i\in E\), are pairwise distinct, and there are only \(2\) residues; by
(T) the residues \(a_i\bmod3\), \(i\in T\), are pairwise distinct, and
there are only \(3\) residues. \emph{Claim 2: \(E\cup T=\{1,2,3,4\}\).}
Given \(i\), pick any \(j\ne i\); by (1) \(g_{ij}\in\{2,3\}\) and
\(g_{ij}\mid m_i\), so \(2\mid m_i\) or \(3\mid m_i\). \emph{Claim 3:
\(|T|\ge2\).} By Claims 1--2, \(|T|\ge4-|E|\ge2\).

We now rule out the remaining sizes of \(T\). - \emph{\(|T|=4\).} Then
all four residues \(a_i\bmod 3\) are pairwise distinct by (T),
impossible with three residues. - \emph{\(|T|=3\).} Say \(T=\{1,2,3\}\)
after relabelling, so \(3\nmid m_4\). For \(j\in T\):
\(g_{4j}\in\{2,3\}\) and \(3\nmid g_{4j}\) (because \(g_{4j}\mid m_4\)),
hence \(g_{4j}=2\), so \(2\mid m_j\). Thus \(2\mid m_1,m_2,m_3\); also
\(2\mid m_4\) by Claim 2 (as \(3\nmid m_4\)). So \(|E|=4\),
contradicting Claim 1. - \emph{\(|T|=2\).} By Claim 2 the two indices
outside \(T\) lie in \(E\), so \(|E|\ge2\), hence \(|E|=2\) by Claim 1
and \(E\) is exactly the complement of \(T\); in particular
\(E\cap T=\emptyset\). Take \(i\in E\) and \(j\in T\). Then
\(g_{ij}\in\{2,3\}\) by (1). But \(j\notin E\) means \(2\nmid m_j\), so
\(2\nmid g_{ij}\) and \(g_{ij}\ne2\); and \(i\notin T\) means
\(3\nmid m_i\), so \(3\nmid g_{ij}\) and \(g_{ij}\ne3\). Contradiction.

Since \(|T|\ge2\) is forced by Claim 3, all cases lead to a
contradiction. Hence some pair satisfies \(g_{ij}\ge4\).
\(\blacksquare\)

\emph{Remark (not needed).} The case \(|T|=3\) shows more generally the
mechanism used: an index whose modulus misses a prime \(p\) forces all
its gcds to avoid \(p\).

\subparagraph{3. Verifica: istruzioni, dipendenze,
tempi}\label{verifica-istruzioni-dipendenze-tempi}

Prova a mano, nessuna dipendenza. Verifica meccanica non necessaria; il
Referee automatico (due giudici indipendenti,
\texttt{runs/p4\_c2/attempts/packet\_001.json}) non ha trovato passaggi
non giustificati.

\subparagraph{4. Fonti e contributo}\label{fonti-e-contributo}

\begin{itemize}
\tightlist
\item
  arXiv:2607.24655 Fornal, Sun, On the problem of large gcd for disjoint
  residue classes (context only: asymptotic result, not used)
\item
  Posizione rispetto allo stato dell'arte: Fornal-Sun {[}2607.24655{]}
  use a gcd-coloured complete graph in the asymptotic regime; the same
  graph viewpoint underlies our finite argument, but the proof here is
  elementary and self-contained.
\item
  Contributo: la prova qui scritta è del team, per esteso.
\end{itemize}

\subparagraph{5. Limiti e parti
irrisolte}\label{limiti-e-parti-irrisolte}

Nessuna: la parte è dimostrata. Gap dichiarati: nessuno.

\subsubsection{Decisioni degli agenti}
\paragraph{Run \texttt{p4\_c2}}
\begin{description}
\item[Iterazione 1 — attempt\_001]
\textbf{Researcher}: famiglia \texttt{direct\_proof}, sotto-obiettivo: Prove the case k=4 by the 'prime cliques' argument: indices whose moduli share a prime p form a set on which residues mod p are pairwise distinct, so |E|\textless{}=2 for p=2 and |T|\textless{}=3 for p=3; a case analysis on |T| closes.. Stato dichiarato: \texttt{CELL\_SOLVED\_CANDIDATE}.
\emph{Perché questa scelta}: With all gcds in \{2,3\}, each modulus is divisible by 2 or 3, and the indices divisible by a common prime p must carry pairwise distinct residues mod p; counting these residues bounds the sizes of the two 'prime cliques' and the remaining case analysis on their sizes is short and exhaustive.
\emph{Rispetto allo stato dell'arte}: Fornal-Sun [2607.24655] use a gcd-coloured complete graph in the asymptotic regime; the same graph viewpoint underlies our finite argument, but the proof here is elementary and self-contained.
\textbf{Referee}: verdetto \texttt{UNKNOWN\_STATUS} (review\_status \texttt{READY\_FOR\_HUMAN}).
\emph{Prossimo passo richiesto}: Human reviews the exact target, proof and evidence, then approves explicit claims
\end{description}

\subsubsection{Codice su cui poggia il tentativo}
Nessun codice nei tentativi.

\section{arXiv literature consulted}
\begin{thebibliography}{99}
\bibitem{2603.26043v1} Yu Hashimoto. \emph{Finiteness of Disjoint Covering Systems with Precisely One Repeated Modulus}. arXiv:2603.26043v1 (2026). \url{http://arxiv.org/abs/2603.26043v1}
\bibitem{1511.04293v1} Shalosh B. Ekhad, Aviezri S. Fraenkel, Doron Zeilberger. \emph{Searching for Disjoint Covering Systems with Precisely One Repeated Modulus}. arXiv:1511.04293v1 (2015). \url{http://arxiv.org/abs/1511.04293v1}
\bibitem{2607.24655v1} Jan Fornal, Yu-Chen Sun. \emph{On the problem of large gcd for disjoint residue classes}. arXiv:2607.24655v1 (2026). \url{http://arxiv.org/abs/2607.24655v1}
\bibitem{2608.15873v1} Murali Menon. \emph{Two Questions on \$G\$-harmonic Tuples}. arXiv:2608.15873v1 (2026). \url{http://arxiv.org/abs/2608.15873v1}
\end{thebibliography}
\end{document}
